\section{Conclusion: every gate becomes software}
\label{sec:conclusion}

The three example routes have different entry points and sequences, and another enterprise would
compose its DGF differently, but each route is still composed of transformations of inputs into evidence, decisions, obligations, and
authorized actions. The gate-substitution conjecture predicts that all of those gates become
machine-executed, including exception handling and support. If that endpoint is attained, the
gates, their sequences, and their purposes can remain while the human labor required to execute
them contracts. This universal endpoint is an argument of the paper, not a measured outcome.

The case for believing it is cumulative. Gate functions already have information contracts; the DGF
combines digital objects, a paved road, decisions about systems, and management pressure as few other
expert functions do; several nearby controls already execute as software; agent horizons are
increasing; class-level authorization can replace repeated individual decisions; and common provision
can fund capabilities that one low-volume enterprise cannot justify alone. None of those facts alone
establishes the endpoint. Together they motivate the comparative hypothesis that an expert
control function will be completely automated here first.

The path is not a straight line. Identical coverage can leave 43.52 or 168.92 workload-equivalent FTE
out of 140, because exceptions, review, rework, and upkeep differ. A residual holding 20\% of the cases
can hold 40\% of the work. The same platform saves on a 95-visit route and loses on a five-visit one.
Separating baseline selection $\kappa$ from post-deployment effort $\eta$ makes those comparisons
coherent, and the account $\phi\eta+(1-\phi)\mu+r+b_A$ names, route by route, the term that still
blocks the endpoint.

The consequence is stated plainly. For a fixed amount of governed work, complete substitution removes
the labor requirement that staffed it: from 140 FTE toward a support floor, and, if support closes,
toward zero. Redeployment, new projects, and changing firm boundaries determine the uses of displaced
capacity; they do not reverse the substitution of the DGF's execution. Governance can even become
stricter while its workforce contracts, because a
control whose marginal cost falls can be applied every time.

The dated hypothesis is an 80\% reduction in required workload-equivalent FTE in the DGF ecosystem
by 20 September 2033 relative to 2026, at comparable governed output, quality, and service. In the
illustrative 140-FTE pool, that means at most 28 FTE, including supervision and support. The existing
configuration figures show how the same labor account can lie above or below this threshold; they do
not date when any configuration will occur. The test measures human hours, converts them consistently
to FTE, and keeps actual staffing and individual headcount as separate outcomes. A ratio above 0.20
rejects the FTE prediction; the current synthetic calculations do not establish its empirical truth.

The companion benchmark provides empirical support for a central technical premise of this
argument: models can successfully perform many of the specified governance reviews. Across 899
evaluable model--case runs, strict gate success ranges from 74.18\% to 94.98\%, while
complete-route success ranges from 24.67\% to 76.92\%. Gemini's 230 completely conforming
routes out of 299 evaluable projects show that its strong performance extends to connected
sequences of reviews, under the stated synthetic conditions. These results support the prospect
of substantial automation of the represented review tasks. The deterministic control also
succeeds on all 1,700 gates and 300 routes, establishing that the supplied rules and facts
already suffice to solve this task. The experiment demonstrates model execution capability,
without identifying an advantage over that direct rule program in this information condition.

The follow-up clarifies both reliability and failure interpretation. A post-hoc structural
evidence criterion raises gate success to 99.06\% for Gemini, 85.53\% for Luna, and 77.24\%
for DeepSeek; these are sensitivity results rather than replacements for the strict scores
or semantic evidence judgments. Authentically observed CSV quotations can fail the original
snapshot-only citation contract. Across 135 repeated runs on 15 existing dossiers, only nine,
three, and zero dossiers, respectively, pass all three strict trajectories. Similar aggregate
performance therefore does not guarantee stable complete success on every project. Conditional
approval and rejected-request counts additionally expose behavior hidden by a final success rate.

The workforce implication is direct but conditional. If comparable performance is achieved in
operational settings, automation could replace review work currently performed by employees
and reduce the staffing required for a comparable volume of governed projects. This makes
workforce contraction a plausible consequence to investigate. Its size depends on the labor
remaining for exceptions, oversight, corrections, integration, and maintenance. Correct
dispositions alone overstate full review conformity: even the best-performing model fails
some evidence checks. No conversion from gate success to an equivalent percentage of FTE
savings is justified by these data.

The benchmark supports the technical prospect of task substitution; the 80\% FTE reduction
by 2033 is the paper's explicit, testable projection of its potential organizational consequence.
Its original 2026 baseline window has closed; testing that exact dated prediction now requires
auditable historical records as well as later quality-constrained human-labor observations.
No such cohort is enrolled. The FTE prediction can succeed while a smaller human team remains indispensable.
Complete execution and support closure remain a longer-term conjecture without a fixed date.
After the DGF, the same reading applies to the professions
of the business itself, from the underwriter to the loan officer (Section~\ref{sec:afterdgf}). The
decisive empirical object is not a model's completion rate. It is the quality-constrained human labor still required by the
connected organizational process; the conjecture predicts its direction, while the labor account
states what must be measured to assess that prediction.
